% Modernised reading — fonds Alexandre Grothendieck, shelfmark 109, pages 1-37
% (the whole folder).
%
% This is an INTERPRETATION, not a transcription. Everything the critical
% apparatus carried moves into footnotes. In any dispute of substance the
% transcription governs: whoever cites Grothendieck must cite batch-01.fr.tex
% and batch-02.fr.tex.
%
% Pass: Opus 5.5 (claude-opus-5-5), 2026-10-03 — first pass, unchecked against the pages by a human.
% The folder was read whole — two batches, pages 1-37, all transcribed; pages
% 11, 15 and 19 are blank separator sheets.
%
% What the folder is. Not a mathematical text but the workshop of an edition:
% the preparation of « Pursuing Stacks » for print, 1983-1984. Pages 1-8 are his
% own working sheets (notations, introduction, printer, index of notations,
% draft titles of §§ 130-133, plan of the volume with two chapter-dependency
% diagrams, questions, terminology). Pages 9-37 come back from correspondents:
% a proposed English version of his presentation text for the series
% « Réflexions Mathématiques » (9-10), a reader's typed comments (12-14), and a
% page-by-page English correction list of the typescript, pp. 1-557, sent in
% instalments dated 19/12/83 and 6/2/84 (16-37). No writer of pages 9-37 is
% named, here or in the transcription. The reading's work is to say, in today's
% words, what the vocabulary and notations of the sheets designate, and to
% check the few mathematical points the correspondents raise.
%
% Notation and conventions fixed for the whole document. « page » is a page of
% the folder; « p. du tapuscrit » a page of the typescript the correction list
% follows; « § » a section of Pursuing Stacks. The typescript is identified
% with Pursuing Stacks as the transcription does (« sans doute »): the leaves
% never name it.
%
% Substantive departures and readings of ours, each footnoted where it occurs:
% — pages 1, 7, 8: the bare numbers of pages 1 and 8 are read as section
%   numbers, from the concordance with page 7, which writes both « p. » and
%   « S. » (S. 35 / aspheric 35; S. 41, 42 / modelizer 28 -> 41, -> 42).
% — page 1: the circle is read as the globe category (the transcription's
%   reading, not written on the sheet).
% — page 3: the underlined N -> double-struck N is read as the typescript's
%   convention for blackboard letters (page 4's doubly underlined Z agrees).
% — page 6: the right-hand dependency diagram runs to IX, the plan to VIII;
%   both stated, neither corrected.
% — page 14: the reader's question whether groupoids form a topos is answered
%   (they do not), with an argument of ours.
% — page 31: « Â is abelian » is false as written (Â is a category of
%   presheaves of sets); the correction to Â_M is endorsed.
% — page 23: j_A(X)(a) = Hom(A_{/a}, X) is stated with its adjunction to the
%   category-of-elements functor, which is ours.
% — page 7: the question « Â cl. mod. cat ? » is said to have been answered
%   later (Cisinski, cited from memory), without implying he knew.
%
% Modern names given and footnoted as ours: test category, weak / local / strict
% test category, modelizer, elementary modelizer, aspheric morphism, homotopy
% interval (after Pursuing Stacks and Maltsiniotis), coherator and
% Grothendieck infinity-groupoid, Thomason structure on Cat, Grothendieck
% fibrations, Dold-Kan, tame topology, anabelian geometry.
%
% Announced and not established: everything the draft titles of §§ 130-133
% promise (page 5); the contents of a « Bibliographie (?) » (page 6); the « ? »
% against « weak equivalence (map of topos) » (page 8); the three questions of
% p. 54 and the question « x |-> Top(B) » (page 7); the formula (6) of typescript
% p. 538, for which the list proposes three readings and settles none (page 36).
\documentclass[11pt,a4paper]{article}
\input{../preamble/grothendieck}

\folder{109}
\pages{1}{37}
\dating{1983-1984}
\watermark{Édition de démonstration\\interprétation personnelle de l'œuvre}
\foldertitle{Tables [Corrections et préparation à l'édition] : notes manuscrites (s.d.), copies de notes manuscrites annotées (s.d.), copies de tapuscrit annoté (s.d.), copies de notes manuscrites (1983-1984) — lecture modernisée du dossier entier}

\begin{document}

\section*{Résumé}

\begin{resume}
En 1983, après douze ans de silence mathématique, Grothendieck écrit en
anglais un long texte, \emph{Pursuing Stacks} (« À la poursuite des champs »),
né d'une lettre à Daniel Quillen et poursuivi au fil des mois. Il pense alors
le publier comme premier volume d'une série, les « Réflexions
Mathématiques ». Ce dossier est l'atelier de cette édition. On n'y démontre
rien : on y prépare un livre.

Un livre de mathématiques ne se réduit pas à ses théorèmes. Il lui faut un
plan, une table des notations, un index des termes, des consignes pour
l'imprimeur, et, quand l'auteur écrit dans une langue qui n'est pas la sienne,
des lecteurs qui relisent la langue. Les huit premières pages sont de sa main
et dressent cet appareil : la liste des signes du livre, des notes pour
l'introduction, un plan du volume avec deux schémas de dépendance des chapitres,
des titres provisoires pour quatre paragraphes, une liste de questions, et
surtout une liste anglaise de termes avec leurs renvois. Cette dernière est
précieuse, parce qu'elle donne en une page le vocabulaire qu'il a forgé pour
le projet : \emph{modélisateur}, \emph{catégorie test}, \emph{asphérique},
\emph{intervalle homotopique}, \emph{cohérateur}.

L'idée qui porte ce vocabulaire se dit sans technique. Les « formes » de la
topologie, à déformation près --- ce qu'on appelle des types d'homotopie ---
se laissent décrire par des assemblages de briques élémentaires : des
triangles et leurs analogues en dimension supérieure, ou des cubes, ou
d'autres formes encore. Grothendieck demande quelles formes de briques
conviennent, et il donne un nom à la propriété qui fait qu'une forme convient
: c'est une \emph{catégorie test}. Une catégorie où l'on peut ainsi décrire
toutes les formes est un \emph{modélisateur}. Les triangles conviennent, les
cubes aussi ; la surprise est que bien d'autres choix conviennent, et que les
petites catégories elles-mêmes, objets purement combinatoires, forment un
modélisateur.

Le reste du dossier vient d'ailleurs. Une version anglaise retouchée du texte
de présentation de la série, qui annonce, au-delà des deux volumes de
\emph{Pursuing Stacks}, deux « continents nouveaux » : une topologie
« modérée », débarrassée des pathologies de la topologie générale, et l'action
des groupes de Galois sur les groupes fondamentaux de certaines courbes. Puis
les commentaires dactylographiés d'un lecteur sur les mérites comparés des
triangles et des cubes. Enfin une longue liste de corrections d'anglais, page
après page sur plus de cinq cents pages de tapuscrit, envoyée en plusieurs fois
entre décembre 1983 et février 1984, où se glissent quelques vraies questions
de mathématiques.

Pour aller plus loin, les noms à chercher sont ceux de la théorie de
l'homotopie de Grothendieck telle qu'elle a été reprise depuis : catégories
test, modélisateurs, et les $\infty$-groupoïdes de Grothendieck.

\keywords{test category, modelizer, aspheric morphism, homotopy interval, coherator, Grothendieck infinity-groupoid, Thomason model structure, cubical set, Dold–Kan correspondence, tame topology, anabelian geometry}
\end{resume}

\section*{Le fil du dossier, et les conventions}

Le dossier compte trente-sept pages, dont trois feuillets de séparation
vierges (pages 11, 15 et 19). Il se partage en deux mains.

\begin{enumerate}
\item[1.] \emph{L'appareil du livre, de sa main} (pages 1 à 8) : notations
(page 1), introduction (page 2), imprimeur (page 3), index des notations
(page 4), titres des §§ 130 à 133 (page 5), plan du volume (page 6), questions
(page 7), terminologie (page 8).
\item[2.] \emph{Le texte de présentation} (pages 9 et 10) : une version
anglaise revue de son texte pour la série « Réflexions Mathématiques »,
recopiée par un correspondant.
\item[3.] \emph{Les retours des lecteurs} (pages 12 à 37) : les commentaires
dactylographiés d'un lecteur (pages 12 à 14), puis une liste manuscrite de
corrections d'anglais suivant le tapuscrit page à page (pages 16 à 37).
\end{enumerate}

Les pages 9 à 37 sont toutes d'autres mains que la sienne, et l'on n'y a
relevé aucune marque de lui. Aucun de leurs auteurs n'est nommé ici. Elles sont
résumées, non reproduites.

\emph{Conventions.} « Page » désigne une page du dossier ; « p.\ du
tapuscrit » une page du tapuscrit anglais que suit la liste de corrections ;
« § » un paragraphe de \emph{Pursuing Stacks}. Le tapuscrit n'est jamais
nommé sur les feuillets ; le groupe de l'inventaire (« Autour de Pursuing
stacks ») et les notions qu'il cite le désignent sans doute
possible.\footnote{La transcription écrit « sans doute » une fois et s'en
tient là ; on fait de même : l'identification repose sur le contenu des renvois
(catégories test, asphéricité, $j_{A}$, intervalles homotopiques), pas sur un
titre.}
Pour une petite catégorie $A$, $\widehat{A}$ est la catégorie des
préfaisceaux d'ensembles sur $A$, $A_{/a}$ la catégorie des objets de $A$
au-dessus de $a$, et $(\mathrm{Hot})$ la catégorie homotopique des types
d'homotopie.

\emph{Annoncé et non établi} : tout ce que promettent les titres des §§ 130 à
133 ; la bibliographie, suivie d'un point d'interrogation ; la notion
d'équivalence faible de topos, laissée sans renvoi ; les questions de la
page 7 ; la formule d'une page tardive du tapuscrit que la liste de
corrections ne parvient pas à lire (page 36).

\pagerange{1}{1}
\section*{Les signes du livre (page 1)}

Une première liste de notations, chacune suivie d'un renvoi :
\begin{itemize}
\item $\Pi_{n}(X)$, et un produit noté $*$ avec indices, sans renvoi ;
\item $(\mathrm{Cat})$, la catégorie des petites catégories, et les
ensembles simpliciaux, au § 16 ;
\item $\Delta$, $\square$, $\bigcirc$, la catégorie des simplexes, celle des
cubes et celle des globes ;\footnote{Les trois signes sont tracés d'un trait
doublé sur un côté, comme des lettres ajourées. L'identification de
$\bigcirc$ à la catégorie des globes --- celle dont les préfaisceaux sont les
ensembles globulaires --- est celle de la transcription ; la feuille ne la dit
pas. Le mot « (Ss-sets) » qui suit $(\mathrm{Cat})$ est lu avec doute.}
\item $A_{/X}$, $A^{\circ}$ (la catégorie opposée), $A^{\wedge} = \widehat{A}$ ;
\item la réalisation géométrique $K \mapsto |K|$, les catégories
$(\mathrm{Ord})$ des ensembles ordonnés et $(\mathrm{Preord})$ des ensembles
préordonnés, $(\mathrm{Hot})$, les simplexes $\Delta_{n}$ et $\Delta[N]$, au
§ 22 ;
\item $\square$, $\square_{n}$, au § 29, et $\tilde{\square}$,
$\tilde{\Delta}$, au § 34.\footnote{Les nombres de la feuille ne disent pas
s'ils renvoient à des pages ou à des paragraphes. On les lit comme des numéros
de paragraphes, parce que la terminologie de la page 8 renvoie aux mêmes
numéros (« standard simplices / cubes », 34 ; « elementary modelizer », 29)
et que la page 7, qui distingue « p. » et « S. », confirme ce sens pour
plusieurs d'entre eux. Cette lecture est la nôtre. Le tilde n'est expliqué
nulle part dans le dossier.}
\end{itemize}

\pagerange{2}{3}
\section*{L'introduction et l'imprimeur (pages 2 et 3)}

Trois notes pour l'introduction : les chiffres soulignés de la table des
matières seront expliqués ; un signe sera défini ; les modifications seront
numérotées et signalées dans une lettre.\footnote{« numérotées » est lu avec
doute ; le destinataire de la lettre n'est pas écrit.} Puis une formule :
\begin{quote}
Stacks = non commutative chain complexes (cf.\ p.\ 4).
\end{quote}
Elle résume le programme du livre. Par la correspondance de Dold et Kan, les
complexes de chaînes de groupes abéliens en degrés positifs équivalent aux
groupes abéliens simpliciaux, c'est-à-dire, du point de vue de l'homotopie,
aux $\infty$-groupoïdes dont la composition est commutative en un sens fort ;
les champs, et plus généralement les $\infty$-groupoïdes, en sont la
généralisation non commutative.\footnote{Ce commentaire est le nôtre. La
formule est sur la page telle quelle, en anglais ; « p.\ 4 » renvoie au
tapuscrit, non à la page 4 du dossier.}

La page 3 donne à l'imprimeur deux correspondances : un $N$ souligné devient
$\mathbb{N}$, et un $\ell$ barré devient un $\ell$ simple. Deux autres signes
--- un double jambage à traits horizontaux, un petit cercle hachuré --- sont
dessinés sans correspondance.\footnote{Le soulignement est la convention du
tapuscrit pour les lettres ajourées, que la machine ne sait pas frapper ; le
$\mathbb{Z}$ doublement souligné de la page 4 va dans le même sens. La
lecture de la convention est la nôtre.}

\pagerange{4}{4}
\section*{L'index des notations (page 4)}

Sur une carte perforée, une liste plus longue, celle d'un index :
\begin{itemize}
\item deux signes d'isomorphie, $\approx$ et $\simeq$, seuls et sur des
flèches ;\footnote{La page ne dit pas lequel désigne l'isomorphisme et lequel
l'équivalence faible ; on ne tranche pas.}
\item $\mathbb{Z}$, $\mathbb{R}$, $\mathbb{Q}$, $\mathbb{C}$ ;
\item $\Delta$, $\square$, $\bigcirc$ et leurs objets $\Delta_{n}$,
$\square_{n}$, $\bigcirc_{n}$ ;
\item $A^{\wedge}$ et un second signe $A^{\&}$, sans
explication\footnote{Le signe est rendu par la transcription comme une
esperluette ; on ne sait pas ce qu'il désigne.} ;
le Hom interne $\underline{\mathrm{Hom}}(A, B)$, pour $A$, $B$ dans
$(\mathrm{Cat})$ ou dans une catégorie quelconque ;
\item les indices bas et hauts $X_{*}$, $X^{*}$, $L_{\bullet}$, $L^{\bullet}$,
c'est-à-dire la graduation homologique et la graduation cohomologique, et le
dual $V^{\vee}$ ;
\item $(\mathrm{Hot})$ et $(\mathrm{Hotab})$, catégorie homotopique des types
d'homotopie et son analogue abélien ; pour une catégorie abélienne $A$, les
catégories dérivées
$D_{\cdot}(A)$, $D^{\cdot}(A)$, $D^{+}(A)$, $D^{-}(A)$, $D(A)$, les
catégories de complexes $\mathrm{Ch}_{\cdot}(A)$, $\mathrm{Ch}^{\cdot}(A)$ et
les catégories homotopiques $K^{+}(A)$, $K^{-}(A)$, $K(A)$, $K_{\cdot}(A)$,
$K^{\cdot}(A)$ ;
\item trois mots anglais : \emph{fibering}, \emph{cofibering},
\emph{bifibering functor} --- les foncteurs fibrés, cofibrés et bifibrés,
qu'on appelle aujourd'hui fibrations et cofibrations de Grothendieck ; et
\emph{localization} « en deux sens », \emph{colocalization}.\footnote{Les deux
sens ne sont pas dits. On peut penser à la localisation d'une catégorie par une
classe de flèches et à la sous-catégorie réflexive (« localisation » au sens
des topos) ; c'est une conjecture, non une lecture.}
\end{itemize}

\pagerange{5}{5}
\section*{Quatre titres de paragraphes (page 5)}

Écrits entre les lignes d'une liste d'ordinateur, des titres anglais pour les
§§ 130 à 133 :
\begin{itemize}
\item[130.] un plaidoyer pour les fibrés non connexes, comme composantes de
modèles semi-simpliciaux ;
\item[131.] un essai de construction d'un « foncteur sphérique »
$\tilde{S}$, et une extension de la notion de fibré ;
\item[132.] un bicomplexe « fou », placé dans le mauvais quadrant, pour les
groupes d'homotopie d'une sphère ;
\item[133.] un titre au nom illisible.\footnote{Le titre de 131 est un
palimpseste dont l'ordre des mots est reconstruit ; un premier état du 132,
annulé par un grand trait, parlait d'une « formule » pour les groupes
d'homotopie des sphères. Le nom du 133 est lu « Saleyman », avec doute ; on
ne l'interprète pas.}
\end{itemize}
Aucun de ces paragraphes n'est dans le dossier. Le 132 vise l'un des
problèmes les plus durs de la topologie, le calcul des groupes d'homotopie
des sphères, et l'adjectif dit assez qu'il s'agit d'une tentative ; on ne
sait pas, par ces feuillets, ce qu'elle contenait.

\pagerange{6}{6}
\section*{Le plan du volume (page 6)}

Le livre devait s'ouvrir sur une dédicace, un sommaire, une table des matières
détaillée et deux introductions, puis le chapitre I, un appendice --- les
lettres à Breen ---, les chapitres II à VIII, des notes et un glossaire, un
index terminologique, un index des notations et une bibliographie suivie d'un
point d'interrogation.\footnote{Une boucle part de « Chap I », entoure
« Appendice : Lettres à Breen » et se ferme sur « Chap II à VIII » : on la lit
comme une insertion de l'appendice entre le chapitre I et les suivants. La
seconde introduction n'a qu'un trait au lieu du mot.}

Au bas, deux schémas de dépendance des chapitres. Dans celui de gauche, en
flèches doubles, II mène à III, III à IV, IV à VII et à VIII, V à VII ; I et
VI ne sont reliés à rien.\footnote{Entre V et VI, un chiffre romain est
lourdement biffé.} Dans celui de droite, en simples traits, la suite I, II,
III, IV est linéaire, IV se relie à VII et à VIII, V à VII, puis VII et VIII
à un chapitre IX ; VI est encore isolé. Les deux schémas ne sont donc pas le
même : le second compte neuf chapitres là où le plan en annonce
huit.\footnote{On rapporte l'écart sans le corriger : rien ne dit lequel des
deux états est le dernier.}

\pagerange{7}{8}
\section*{Questions et terminologie (pages 7 et 8)}

\subsection*{Quatre questions}

La page 7 relève, par page du tapuscrit et par paragraphe, des points à
revoir :
\begin{itemize}
\item p.\ 54, § 35 : trois questions sur cette page, dit-il, sans les
écrire ;
\item p.\ 55, § 35 : $x \mapsto \mathrm{Top}(B)$, sans plus ;\footnote{On ne
sait pas quelle construction la formule abrège ; on ne la reconstruit pas.}
\item p.\ 72, § 41 : $\widehat{A}$ est-il une catégorie de modèles fermée ?
\item p.\ 73--74, § 42 : $(\mathrm{Ord})$ est-il un modélisateur ?
\end{itemize}

La troisième est la question de savoir si, pour une catégorie test $A$, les
équivalences faibles de $\widehat{A}$ s'insèrent dans une structure de
modèles de Quillen. Elle a reçu depuis une réponse positive, avec les
monomorphismes pour cofibrations.\footnote{D.-C.~Cisinski, \emph{Les
préfaisceaux comme modèles des types d'homotopie}, Astérisque 308 (2006),
cité de mémoire. La réponse est postérieure de vingt ans au dossier, et rien
n'indique qu'il l'ait connue ; « cl.\ mod.\ cat » est l'abréviation de la
\emph{closed model category} de Quillen.} La quatrième demande si les ensembles
ordonnés, avec pour équivalences faibles les applications dont le nerf est une
équivalence faible, suffisent à décrire tous les types d'homotopie ; la
réponse est oui, parce que la subdivision barycentrique d'un ensemble
simplicial a même type d'homotopie qu'un nerf d'ensemble ordonné.\footnote{La
réponse est la nôtre, sous cette forme rapide ; la feuille pose la question
sans y répondre, et l'on ne sait pas par elle ce qu'en dit le § 42.}

\pagerange{8}{8}
\subsection*{Le vocabulaire du livre}

La page 8, « Terminology », dresse en anglais la liste des termes du livre,
dans l'ordre de leur apparition, avec le paragraphe où chacun est
introduit.\footnote{Les numéros sont dans une colonne à droite, et leur
rattachement aux lignes est donné au plus près. On les lit comme des numéros de
paragraphes pour la raison dite à la page 1 ; les numéros 41 et 42, en
italique sur la ligne de « modelizer », sont ceux des paragraphes de la
page 7, et « aspheric », 35, celui de la question de la p.\ 54.} On la
regroupe ici par notion, en disant en termes actuels ce que chacune désigne.

\emph{Structures et champs} (§§ 4 à 13). Les structures « primitive,
primaire, secondaire, ternaire… », les types standard (produits fibrés
itérés), les $n$-champs, les $\infty$-champs et les $n$-champs relatifs à une
catégorie $C$. Puis le \emph{cohérateur} (§§ 13 et 25) : la donnée qui dit
quelles opérations et quelles cohérences un $\infty$-groupoïde doit porter.
C'est d'elle que G.~Maltsiniotis a tiré, dans les années 2000, une définition
des $\infty$-groupoïdes dits aujourd'hui « de Grothendieck ».\footnote{Le
renvoi est le nôtre et cité de mémoire ; la feuille ne porte que le mot et deux
numéros.}

\emph{Modélisateurs} (§§ 26 à 29). Un \emph{modélisateur} est une catégorie
$M$ munie d'une classe $W$ de flèches, dites équivalences faibles, telle que
la localisée $M[W^{-1}]$ soit équivalente à $(\mathrm{Hot})$, de façon
compatible aux foncteurs naturels ; un foncteur entre modélisateurs qui
respecte cette structure est dit \emph{modélisant}. Le modélisateur de base est
$(\mathrm{Cat})$, les équivalences faibles étant les foncteurs dont le nerf
est une équivalence faible d'ensembles simpliciaux.\footnote{L'équivalence
$(\mathrm{Cat})[W^{-1}] \simeq (\mathrm{Hot})$ remonte à Quillen et Illusie ;
la structure de modèles correspondante sur $(\mathrm{Cat})$ est due à
R.~W.~Thomason (1980). Références citées de mémoire ; la feuille ne nomme
personne.} Une \emph{catégorie test} (§ 27) est une petite catégorie $A$
telle que $\widehat{A}$, muni des flèches que le foncteur « catégorie des
éléments » $i_{A} : \widehat{A} \to (\mathrm{Cat})$ envoie sur des
équivalences faibles, soit un modélisateur, cela restant vrai pour chaque
$A_{/a}$, et $A$ ayant un nerf contractile. Un \emph{modélisateur
élémentaire} (§ 29) est un tel $\widehat{A}$. La ligne de la feuille garde
la trace d'une hésitation : « test-topology », puis « test category », suivi
d'un point d'exclamation.\footnote{Les définitions sont dites dans la forme
qu'elles ont prise dans \emph{Pursuing Stacks} et dans l'exposé de
G.~Maltsiniotis, \emph{La théorie de l'homotopie de Grothendieck}, Astérisque
301 (2005), cité de mémoire : catégorie test faible quand seule la condition
sur $A$ est demandée, catégorie test locale quand elle l'est pour chaque
$A_{/a}$, catégorie test quand $A$ est de plus asphérique. La feuille ne donne
que les noms.}

\emph{Catégories test strictes, simplexes et cubes} (§§ 33 et 34). Une
catégorie test est \emph{stricte} quand la localisation
$\widehat{A} \to (\mathrm{Hot})$ commute aux produits finis. Les
simplexes et les cubes standard sont les exemples de base : $\Delta$ est une
catégorie test stricte, et la catégorie des cubes $\square$ est une catégorie
test qui n'est pas stricte, le produit des ensembles cubiques n'ayant pas le
bon type d'homotopie.\footnote{L'énoncé sur $\square$ est le nôtre et suit
la littérature ; il rejoint la remarque du lecteur de la page 13. Deux carrés
biffés suivent la ligne, au même numéro 34.}

\emph{Asphéricité dans les topos} (§§ 35 à 39). Un topos est
\emph{asphérique} quand il a le type d'homotopie d'un point ; un objet d'un
topos l'est quand le topos induit au-dessus de lui l'est, et une flèche d'objets
quand ses fibres le sont, au sens convenable. Suivent les topos localement
asphériques et totalement asphériques, puis les \emph{équivalences faibles}
de topos, sans renvoi : la ligne porte un point d'interrogation à la place du
numéro. L'\emph{intervalle homotopique} (§ 37) est un objet $I$ de
$\widehat{A}$ muni de deux sections disjointes $0, 1$ de l'objet final,
asphérique, qui joue le rôle du segment $[0,1]$ et permet de parler
d'homotopies dans $\widehat{A}$ ; dans les ensembles simpliciaux,
c'est $\Delta_{1}$.\footnote{La définition est donnée dans sa forme reçue ;
la feuille ne porte que le terme et le numéro 37.}

\emph{Variantes locales} (§ 39). Les modélisateurs élémentaires stricts, les
catégories test locales et les modélisateurs élémentaires locaux, les topos
modélisants et localement modélisants. Le § 39 apparaît aussi en italique
sur la ligne des « Hot-équivalences de structures de site » (§ 26), comme un
renvoi complémentaire.

\pagerange{9}{10}
\section*{Le texte de présentation (pages 9 et 10)}

Deux pages d'une autre main, qui lui sont adressées : le correspondant
recopie une version anglaise de son texte de présentation, la photocopie de
l'original étant devenue peu lisible, et souligne en rouge les mots changés
par rapport à cet original.\footnote{La note d'en-tête dit que les changements
viennent de plusieurs relecteurs, sans les séparer sur la page. Les
retouches sont pour l'essentiel de langue. Les mots soulignés en rouge sont en
italique dans la transcription ; le texte, de cette main, n'est pas reproduit
ici.} Le milieu manque : une ligne dit que le texte continue sans changement
jusqu'au dernier paragraphe.

Le texte présente l'œuvre de la maturité d'un auteur qui revient après douze
ans de silence, comme le journal, au jour le jour, d'un voyage d'exploration
mêlé de réflexions sur le voyage lui-même. Les deux volumes à venir, sous le
titre commun \emph{Pursuing Stacks}, le premier sous-titré \emph{The
Modelizing Story}, esquissent une synthèse de l'algèbre homotopique, de
l'algèbre homologique et de la théorie des topos, dont le fil conducteur vient
de la géométrie algébrique et des notions de faisceau, de champ, de 2-champ.

Ces volumes ouvrent une série plus large, les « Réflexions Mathématiques »,
qui doit présenter des idées sur deux « continents nouveaux ». Le premier est
une notion d'espace topologique « modéré » (\emph{tame}), cadre naturel de
l'intuition de la forme, que la topologie générale, faite pour l'analyse,
encombre de pathologies ; elle servirait aux espaces d'homéomorphismes et de
plongements et à une théorie du dévissage des structures stratifiées, qu'on
rencontre partout où il y a des modules. Le second est l'étude de l'action
des groupes de Galois profinis absolus, comme
$\mathrm{Gal}(\overline{\mathbb{Q}}/\mathbb{Q})$, sur certains groupes
fondamentaux géométriques profinis dits anabéliens, comme celui de
$\mathbb{P}^{1}_{\overline{\mathbb{Q}}} \smallsetminus \{0, 1, \infty\}$.
Le dernier paragraphe explique la langue : les deux premiers volumes, nés
d'une correspondance en anglais, sont écrits en anglais, mais la première
partie de l'introduction, de caractère personnel, l'est en
français.\footnote{Les deux continents sont ceux de l'\emph{Esquisse d'un
programme}, de 1984 : la topologie modérée, dont les structures o-minimales
sont aujourd'hui une forme, et la géométrie anabélienne, avec les dessins
d'enfants et la tour de Teichmüller. Ces rapprochements sont les
nôtres. La série « Réflexions Mathématiques » n'a pas paru sous ce nom ;
\emph{Pursuing Stacks} a été publié bien plus tard par la Société
mathématique de France, information citée de mémoire.}

\pagerange{12}{14}
\section*{Les commentaires d'un lecteur (pages 12 à 14)}

Trois pages dactylographiées, « Comments on "Pursuing Stacks" », d'un lecteur
qui n'est pas nommé. Après une liste de fautes d'orthographe, quatre remarques
de fond.

\emph{Ensembles simpliciaux ou cubiques.} Les ensembles simpliciaux doivent
leur place à des propriétés précises : leur catégorie est cartésiennement
fermée ; la réalisation géométrique induit une équivalence entre la catégorie
homotopique des complexes de Kan et celle des CW-complexes, et commute aux
produits finis quand on prend la topologie compactement engendrée ; il y a
une correspondance de Dold--Kan et une démonstration simpliciale de la suite
spectrale de Serre. Pour les ensembles cubiques, l'équivalence homotopique
n'était alors démontrée en détail que dans un mémoire non publié ; le produit
cartésien n'a pas la bonne réalisation, seul un produit tensoriel l'a ; et une
correspondance de Dold--Kan demande la structure supplémentaire des
« connexions » de Brown et Higgins.\footnote{Ces énoncés sont exacts tels
qu'ils sont rapportés. Le lecteur rappelle aussi la description de Kan du
foncteur des complexes de chaînes vers les groupes abéliens simpliciaux,
$C \mapsto \bigl(n \mapsto \mathrm{Hom}(N\mathbb{Z}\Delta^{n}, C)\bigr)$,
avec $N$ les chaînes normalisées ; la page écrit $\mathrm{Hom}(C_{N}\Delta^{n},
C)$. Il préférerait qu'on parle du théorème de Dold--Kan plutôt que de Puppe.}

\emph{Modèles algébriques des types d'homotopie.} Il s'étonne que l'approche de
Loday passe par des objets $n$-simpliciaux, rappelle que la théorie de van
Kampen de Brown et Loday emploie des groupoïdes $n$-uples, les
$\mathrm{Cat}^{n}$-groupes, et se demande si l'extension de tels objets aux
topos ne demanderait pas une construction d'intervalle unité, là où les
champs entreraient en jeu.

\emph{Le nerf.} Prendre le nerf de catégories qui modélisent une situation
géométrique a servi en topologie géométrique, pour la K-théorie algébrique des
espaces, sans être devenu un outil de calcul. Il propose un test : calculer
$\pi_{3}(S^{2})$ sans connaître la fibration de Hopf $S^{1} \to S^{3} \to
S^{2}$.\footnote{On sait que $\pi_{3}(S^{2}) \simeq \mathbb{Z}$, engendré
par la fibration de Hopf ; le test est donc d'y arriver par une autre voie.}

\emph{Les groupoïdes et les topos.} À propos d'une page de l'introduction, il
demande si la catégorie des groupoïdes n'est pas un topos faute de
classifiant des sous-objets. Elle n'en est pas un, et c'est bien l'absence de
classifiant qui l'empêche.\footnote{L'argument est le nôtre. Les sous-objets du
groupoïde ponctuel sont deux, donc un classifiant $\Omega$ aurait deux objets.
Le groupe $\mathbb{Z}/3$, vu comme groupoïde à un objet, a trois sous-objets :
le vide, l'objet seul et le tout. Une flèche de lui vers $\Omega$ est le choix
d'un objet $x$ et d'un $h \in \mathrm{Aut}(x)$ avec $h^{3} = 1$. Chaque objet
en fournit au moins une, l'identité, et il en faudrait trois en tout : l'un des
deux objets devrait avoir exactement deux solutions de $h^{3} = 1$. C'est
impossible : les solutions autres que $1$ vont par paires $\{h, h^{2}\}$, avec
$h^{2} \neq h$, et leur nombre est impair.}
Au pied de la dernière page, d'une main qui ne semble pas la sienne : « so
far, so good! »

\pagerange{16}{37}
\section*{La liste de corrections (pages 16 à 37)}

\subsection*{La liste et ses envois}

D'une même main que les pages 9 et 10, une liste en anglais de corrections au
tapuscrit, page par page et ligne par ligne, adressée à lui. Elle commence
par « Errors in 1st 100 pages of notes + letter to Quillen » et suit le
tapuscrit sans lacune de la p.\ 1 à la p.\ 557, où elle s'arrête au tiers d'une
feuille. Elle est envoyée en plusieurs fois : la page 20 porte « Sent
19/12/83 » après la p.\ 100, la page 28 « Sent 6/2/84 » après la
p.\ 354.\footnote{La page 20 est la copie d'un état antérieur des entrées des
p.\ 87 à 100, annulé de deux traits, suivi de la suite jusqu'à la p.\ 115 ; la
page 27 est de même la copie d'un état antérieur des vingt premières lignes de
la page 28.} Neuf entrées sur dix sont d'usage : « commute with » et non
« commute to », « standard » et non « standart », « aspheric » et non
« asperic » (« avec les asperges ! », ajoute le correspondant en français),
les contractions, « formulary » pour « formulaire », trouvé excellent. Les
autres touchent aux mathématiques ; on les rassemble par thème.

\pagerange{21}{26}
\subsection*{Catégories test et asphéricité (pages 21 à 26)}

\emph{Le foncteur $j_{A}$} (p.\ 224 du tapuscrit). Le correspondant rappelle,
en renvoyant à la p.\ 32,
\[
j_{A} : (\mathrm{Cat}) \longrightarrow \widehat{A}, \qquad
j_{A}(X)(a) = \mathrm{Hom}_{(\mathrm{Cat})}(A_{/a}, X),
\]
et juge fausse en conséquence une ligne du tapuscrit. Ce $j_{A}$ est l'adjoint
à droite du foncteur $i_{A} : \widehat{A} \to (\mathrm{Cat})$ qui associe à
un préfaisceau sa catégorie des éléments : $i_{A}$ prolonge par colimites le
foncteur $a \mapsto A_{/a}$, et $j_{A}$ en est l'adjoint par
construction.\footnote{L'adjonction est notre commentaire ; elle est au
centre de la théorie des catégories test, où l'on demande que $i_{A}$ induise
une équivalence des localisées. Quelques mots interlinéaires sous la formule
sont illisibles.}

\emph{Asphéricité} (p.\ 234 et 235). Une reformulation est proposée d'un
énoncé sur $W$ et un couple $(A, i)$ ; l'énoncé « $x$ asphérique si et
seulement si $a$ asphérique » est relevé comme une tautologie telle qu'il est
écrit.\footnote{Sans le contexte du tapuscrit, on ne peut dire quelle
condition était voulue ; la remarque signale une coquille, non une erreur de
fond.}

\emph{Équivalences faibles induites} (p.\ 240 à 251). La formule
$W = (i_{M,i})^{-1}(\underline{W})$, qui définit les équivalences faibles
d'une catégorie par image réciproque de celles d'une autre ; une question
sur la condition (ii) d'une définition, qui demanderait $i_{0} : N_{0} \to N$ ;
$i : A \to M$ à la p.\ 251.

\emph{Notations à fixer.} $\pi_{0}$ et $\pi_{1}$ ne sont pas définis à
l'endroit où ils servent (p.\ 269) ; faut-il lire $\alpha : \Delta \to
\Delta_{/F}$ (p.\ 327) ? Une formule $K_{h} \subset Z_{h}$ est à corriger
(p.\ 116)\footnote{L'indice de gauche est lu avec doute ; la transcription
lit $K_{h} \subset K_{n}$ sur la page, avec la correction proposée en $Z_{h}$.} ;
une autre, $f^{\circ} : C^{\circ} \to C'^{\circ}$, récrite (p.\ 167) ; une
définition des intervalles homotopiques non faibles est à tourner autrement
(p.\ 186) ; un renvoi du tapuscrit au § 51 viserait le § 52 (p.\ 114--115).
En tête de la page 25, une relation « $C \simeq A$ ? » est jugée
confuse.\footnote{Le signe de relation est surchargé ; la lecture est
incertaine.}

\pagerange{29}{32}
\subsection*{Algèbre homotopique abélienne (pages 29 à 32)}

\emph{Préfaisceaux abéliens} (p.\ 402 et 428). L'indice de
$\widehat{A}_{\mathrm{ab}}$ manque sur la copie ; et « $\widehat{A}$ is
abelian » est une erreur, c'est $\widehat{A}_{M}$ qui est voulu. La
correction est juste : $\widehat{A}$, catégorie de préfaisceaux d'ensembles,
n'est pas même additive ; ce sont les préfaisceaux en modules qui forment une
catégorie abélienne.\footnote{L'indice $M$ désigne vraisemblablement un
anneau ou un module de coefficients ; la liste ne le dit pas.}

\emph{Formules} (p.\ 385, 413, 417, 438). Un signe $\simeq$ manque dans la
formule (7), une parenthèse ailleurs, et les renvois aux formules (75) et
(84) sont relevés.

\emph{Groupes d'homotopie et objets en groupes} (p.\ 449 et 471). Une phrase
sur des groupes d'homotopie qui ne sont pas des $k$-modules est reformulée.
Une autre, jugée ambiguë, est récrite ainsi : un complexe $X_{*}$, lisse et
satisfaisant à la condition de Kan, devrait donner un objet en groupes
$G = \pi_{1}(X_{*})$ dans $U(k)$, avec une submersion notée (H) ; le
correspondant ajoute qu'il a pu se méprendre.\footnote{On ne sait pas, par
le dossier, ce que désigne $U(k)$ ; on rapporte la reformulation sans la
juger.}

\pagerange{33}{37}
\subsection*{Puissances divisées, et une formule illisible (pages 33 à 37)}

\emph{Puissances symétriques et divisées} (p.\ 484 à 494). Une flèche manque
dans la formule (8) ; une expression contenant $\Gamma^{\wedge}_{k}(M)$ est
relevée (p.\ 485)\footnote{La transcription lit
$W^{\wedge}\,\Gamma^{\wedge}_{k}(M)$, avec doute.} ; une parenthèse est mal
placée dans $\prod_{j} \Gamma^{j}_{k}(M)$. Surtout, le correspondant demande
qu'on explique les notations $\mathrm{Sym}^{j}_{k}$ et $\Gamma^{j}_{k}$, dont
il a dû chercher le sens ailleurs. Ce sont les puissances symétriques et les
puissances divisées d'un $k$-module, $\Gamma_{k}(M) = \bigoplus_{j}
\Gamma^{j}_{k}(M)$, et le chapeau désigne vraisemblablement le produit
complété $\prod_{j}$.\footnote{La liste renvoie à SGA 6 pour $\Gamma_{k}$.
L'interprétation du chapeau est la nôtre, tirée du $\prod_{j}$ de la même
entrée.}

\emph{Schémas en groupes et dévissages} (p.\ 495 à 517). Le mot
« endomorphisme » est mis en doute, les groupes d'endomorphismes étant rares,
le correspondant avouant mal connaître les schémas en groupes additifs ; un
« 0 » manque dans une cohomologie ; « Lazare » ou « Lazard » ? ; la formule (8)
se factorise par (2$'$) et non par (2) ; que désigne $V'$, quand $V$ et $V'$
servent tous deux ?

\emph{Une condition d'annulation} (p.\ 538). La formule (6) paraît étrange au
correspondant, l'indice de $\Pi$ étant placé ailleurs qu'il ne l'attend. Il
propose trois lectures sans trancher :
\[
H^{i}\bigl(T, \Pi_{i}(X_{*})\bigr) = 0 \ (i > 0), \qquad
H^{i}\bigl(T, \Pi_{j}(X_{*})\bigr) = 0 \ (j > 0), \qquad
H^{i}\bigl(T, \Pi_{j}(X_{*})\bigr) = 0 \ (i \geq j > 0).
\]
Le dossier ne dit pas laquelle était voulue, et l'on ne choisit
pas.\footnote{La deuxième lecture demande l'annulation de toute la
cohomologie, la troisième une annulation « au-dessus de la diagonale », du
type de celles qui rendent une obstruction nulle dans une tour de Postnikov.
Ce commentaire est le nôtre et ne tranche rien.}

La liste s'achève à la p.\ 557, sur des corrections d'usage dans un passage
de récit et l'orthographe de « Hurewicz ».

\end{document}
